\section{Comparison of models and conclusions}\label{sec:comparison}
\subsection{Probabilistic read-once branching programs}
After fixing the seed and padding layers to a common width, the product
in Definition~\ref{def:machine} is a stochastic ordered read-once
branching program: layer $t$ reads command $c_t$ once and applies
$T_{t-1,c_t}$. This is the randomized-transition extension of the
standard nonuniform model in \cite[Section 1.1]{RSV}; see also
\cite{Wegener}. Width is the maximum layer size, whereas total size
counts all layers. Our theorem is a width theorem, not a succinct
uniform program-size or bit-operation theorem.

For a chosen query $j$, appending an accepting Bernoulli row with
probability $(1+d_j(s))/2$ turns the real readout into ordinary randomized
acceptance with no larger width than $\max(K,2)$. Combining the two
means into $d_1+i d_2$ does not enlarge the retained alphabet. This
combination is an identity between mean functions, not a requirement to
sample both query answers jointly. Thus real readouts and the word
``transducer'' do not by themselves separate our problem from branching
programs.

For a fixed initial row $p$ and terminal column $d$, put
$F(c)=pT_{0,c_1}\cdots T_{N-1,c_N}d$. Its Boolean-input Fourier
coefficient at $A\subseteq[N]$ is exactly
\begin{equation}\label{eq:bp-fourier}
 \widehat F(A)=p\prod_{t=1}^N
 \frac{T_{t-1,0}+(-1)^{\one_{\{t\in A\}}}T_{t-1,1}}2\,d.
\end{equation}
This follows by multiplying independent sums over input bits in their
original order. It is a matrix-product Fourier coefficient, distinct
from a harmonic of the conditional phase measure used in \cite{GTF34}.
The present proof instead uses the nonlinear quantity
$\E|\E[\zeta^{L_t}\mid S_t]|$ at packet boundaries.

Reingold--Steinke--Vadhan \cite[Theorem 1.4]{RSV} bound Boolean
level-$l$ Fourier mass by $(2w^2)^l$ for regular width-$w$ programs.
Steinke--Vadhan--Wan \cite{SVW} establish width-three Fourier-growth
bounds and use them for pseudorandomness. Lee--Pyne--Vadhan
\cite[Definitions 3--5 and Theorem 6]{LPV} refine the regular-program
bounds, including unbounded width and restrictions on accepting states;
their Section 1.2.3 discusses generalized group products. These results
concern Fourier growth or fooling classes of computations. They do not
state the terminal-correlation packet inequality \eqref{eq:packet} or
its nonuniform occupation consequence.

Regularity means preservation of a uniform layer distribution under a
fair next input. It is not implied by row stochasticity: the matrix all
of whose rows put mass one on a single label is legal here and destroys
uniformity. Consequently a regular/permutation theorem cannot be invoked
without its hypothesis. Conversely, some bounds for general deterministic
programs extend to stochastic programs by resolving all private rows
into a convex combination of deterministic tables and using convexity.
We therefore make no claim that private randomness alone prevents use
of the earlier literature. The extra statement proved here is the
multiplicative loss forced by a separated command packet, for arbitrary
starting conditional centroids and arbitrary endpoint assignments.
It gives the matching growing-width law without a regularity assumption
and without an estimate of the Fourier mass in \eqref{eq:bp-fourier}.

\subsection{Circle walks and the choice of test law}
For independent Bernoulli commands of bias $p$, the uncompressed phase
has harmonic multiplier
\[
 a_l=1-p+p e^{2\pi il\alpha},\qquad
 |a_l|^2=1-4p(1-p)\sin^2(\pi l\alpha).
\]
These are classical circle-walk quantities. Berkes--Borda \cite{BB}
relate convergence and limit theorems for circle walks to rational
approximation, including the transition along convergent denominators.
That analysis concerns the uncompressed walk, not the information
retained by an additional hidden register.

The earlier Fourier budget \cite{GTF34} combined this damping, common
conditional-amplitude loss, and a positive finite-atom Fourier witness.
It gave a fifth-root unrestricted converse. Our test law is different:
its bits are correlated inside a packet so that the packet count is
uniform, and different packets are independent. We are not asserting a
cube-root occupation theorem for the fair-command law itself. Wordwise
correctness allows the new law; Lemma~\ref{lem:packet} measures its
compression directly. The packet costs $O(k)$ commands and loses
$\Omega(k^{-2})$ amplitude, explaining the cube-root order.
The angular averaging estimate remains valid with its sharp universal
coefficient, proved in Appendix~\ref{sec:angular}; improving that
coefficient's quadratic order is not needed for the new result.

\subsection{Positive realization and experiment comparison}
Invariant convex sets, positive stochastic realizations, and irrational
rotation obstructions for one repeated positive matrix belong to
classical realization theory \cite{Heller,Vidyasagar,BF,MW}.
The present lower bound permits unrelated matrices at every layer and a
new machine for each horizon; it is not a peripheral-eigenvalue argument
for a stationary realization. The upper construction is inherited from
\cite{GTF33,GTF34} and has been proved again in
Proposition~\ref{prop:upper}. What is new in the main theorem is its
match against all hidden states, including under fixed row error.

Finite Blackwell garbling and one-sided deficiency are classical
\cite{Blackwell,Torgersen}. Theorem~\ref{thm:final-error} bounds the
actual terminal optimization. Corollary~\ref{cor:diagram} then bounds
an accumulated local certificate; it does not identify that certificate
with the optimized terminal distance. Orbit quantization, sphere-cap
bounds and polytope nets in Section~\ref{sec:compact} are classical
geometric mechanisms. Their representation-valued sequential loss is
stated and proved explicitly in Theorem~\ref{thm:group-budget}.

\subsection{Conclusions and precise boundaries}
The hidden and vector-state rotation widths have the same cube-root
order for each badly approximable angle and fixed nonzero signal at
sufficiently small fixed error. The finite constants are proof constants,
not optimized thresholds or exact integer minima. The packet budget
also controls selected low-width endpoints and actual terminal error.
For compact rotation commands, a distinct matched sphere-dimension law
follows from the same contraction theorem.

No sharp order is asserted for every irrational angle, for every finite
noncommuting command set, or for a regime in which signal and permitted
error vary with the horizon. Optimizing informative action choice is
not the external-command problem studied here. Atomic computable-real
rows are not finite fair-bit algorithms. No exact minimum dynamic
Blackwell distance or independent analytic foundation for an unrelated
model is inferred from these results. The original manuscript and its
historical arguments remain separately available, with their hypotheses
unchanged.
